\documentclass[10pt,twocolumn]{article}

% ---------------------------------------------------------------------------
% Portable wearable baselines on a public ledger.
% Self-contained: uses only base LaTeX packages so it compiles anywhere.
% ---------------------------------------------------------------------------
\usepackage[letterpaper,margin=0.75in]{geometry}
\usepackage[T1]{fontenc}
\usepackage{lmodern}
\usepackage{amsmath,amssymb}
\usepackage{booktabs}
\usepackage{xcolor}
\usepackage{listings}
\usepackage{microtype}
\usepackage[hidelinks]{hyperref}
\usepackage{enumitem}
\setlist{nosep,leftmargin=1.2em}

\definecolor{codebg}{RGB}{245,245,247}
\definecolor{codekw}{RGB}{120,60,200}
\lstset{
  basicstyle=\ttfamily\footnotesize,
  backgroundcolor=\color{codebg},
  breaklines=true,
  columns=fullflexible,
  keepspaces=true,
  frame=none,
  aboveskip=6pt, belowskip=6pt,
  xleftmargin=4pt, xrightmargin=4pt,
}

\newcommand{\code}[1]{\texttt{\small #1}}

\title{\vspace{-2.5em}\bfseries Wearable Baselines on a Public Ledger:\\
Portable, Owner-Controlled Calibration for\\Cross-App BCI and Biosignal Interfaces}

\author{Alexandru Vicol\\
\small University of Toronto, Department of Electrical \& Computer Engineering\\
\small \texttt{HealthLog Protocol} --- built on Solana + RPC Fast\\
\small \emph{Draft, \today}}

\date{}

\begin{document}
\maketitle

\begin{abstract}
\noindent
Every consumer neurotechnology and biosignal application re-derives the same
thing on first launch: a \emph{per-user baseline}. An EEG brain--computer
interface (BCI) fits a personal arousal or intent model; an EMG interface
records a maximum-voluntary-contraction (MVC) trial to normalize effort; an HRV
app spends days establishing resting norms. These baselines are expensive to
collect, physiologically slow to change, and today are trapped inside whichever
app produced them. A user who installs a second app starts from zero. There is
no portable, verifiable, owner-controlled representation of ``my calibration.''

We present a design and working implementation that treats a baseline as a
first-class, portable asset. The fitted parameters stay encrypted on the user's
device; only a SHA-256 \emph{commitment} plus non-identifying shape metadata is
anchored on a public ledger (Solana). Any third-party app that later receives
the baseline out of band re-hashes it and checks it against the on-chain
commitment before trusting it---obtaining integrity, provenance, ownership, and
revocation without the raw biometric ever touching the chain. We extend the
HealthLog consent protocol (Solana Colosseum, 2026) from \emph{sessions} to
\emph{baselines}, using RPC Fast as the mainnet verification rail. We anchor and
round-trip-verify real EEG, EMG, and HRV baselines on Solana devnet, show that
tampering with a single parameter is detected, measure a median verification
read latency of $\approx\!106$\,ms, and enforce the privacy invariant with a
runtime gate and automated tests. We close with an honest account of what the
current commitment-only design does \emph{not} yet provide---zero-knowledge
proofs of baseline quality and cryptographic biometric-template protection---and
a roadmap toward them.
\end{abstract}

\section{Introduction}

Consumer wearables and BCIs have made biosignals---EEG, EMG, HRV, EDA,
PPG---cheap to record. What has \emph{not} become cheap is \emph{calibration}.
Because physiology varies across people and electrode placements, almost every
serious application begins by estimating a user-specific baseline:

\begin{itemize}
\item An EEG focus/relaxation BCI fits a personal model over band-power
  features (e.g.\ a Gaussian mixture over $\theta,\alpha,\beta$ and the
  theta/beta ratio) so that ``high arousal'' means high \emph{for this user}.
\item An EMG gesture or prosthetic interface records an MVC trial and resting
  RMS per channel, so effort can be expressed as a percentage of that user's
  own maximum.
\item An HRV stress/recovery app accumulates a multi-day resting baseline
  (SDNN, RMSSD, LF/HF) so a live reading is scored relative to the user rather
  than to a population.
\end{itemize}

These baselines share three properties. They are \textbf{costly}: minutes of
guided recording for EEG/EMG, days for HRV. They are \textbf{slow-changing}:
useful for weeks to months. And they are \textbf{siloed}: each app stores its
baseline privately, so the cost is paid again for every new app, device, or
research study. A user who has already recorded a perfectly good EEG arousal
baseline in app~A cannot hand it to app~B; app~B re-records from scratch, and
neither app can prove to a third party (a clinician, a research cohort, a second
device) that the baseline it holds is authentic, unmodified, and owned by the
person presenting it.

We argue that a baseline should be a \emph{portable, verifiable, owner-controlled
object}, and that a public ledger is the right substrate for its
\emph{integrity and consent envelope}---though emphatically not for the biometric
data itself. Our contributions:

\begin{enumerate}
\item A \textbf{baseline-commitment design} in which only a SHA-256 hash of the
  canonicalized parameters plus shape metadata (signal type, model family,
  feature-schema id, dimensionality, recalibration horizon) is anchored
  on-chain. The fitted parameters never appear on-chain in any form.
\item An \textbf{integration into HealthLog}, extending its Colosseum consent
  protocol from sessions to baselines, with \textbf{RPC Fast} as the mainnet
  write/verification rail and a free devnet demo rail.
\item A \textbf{working reference implementation and evaluation}: real EEG, EMG,
  and HRV baselines anchored and verified on Solana devnet; single-parameter
  tamper detection; a runtime privacy gate with automated tests; and read-latency
  measurements against RPC Fast.
\item An \textbf{honest security analysis} that names the residual risks---
  biometric-template leakage under parameter sharing, linkage across
  commitments, and the absence of zero-knowledge quality proofs---and a roadmap.
\end{enumerate}

\section{Background and Motivation}

\paragraph{What a baseline is, concretely.} In the HealthLog EEG pipeline, the
baseline is a serialized \code{GmmModelState}: the feature columns, a
standardization (mean/scale per feature), Gaussian-mixture weights, means, and
covariances, and a mapping from cluster index to an arousal label. This is
exactly the object a downstream inference step loads to classify live EEG
without re-training. An EMG baseline is a small vector of per-channel MVC and
resting-RMS values plus the filter settings used to obtain them; an HRV baseline
is a handful of resting time- and frequency-domain statistics with the window
over which they were computed. In all three cases the baseline is (i) small
(hundreds of bytes to a few kilobytes), (ii) sufficient to make a second app's
inference work, and (iii) sensitive, because it is derived from a person's
physiology.

\paragraph{Why not just copy the file?} A user could in principle export the
baseline file and import it elsewhere. Two things are missing. First,
\emph{integrity and provenance}: the receiving app has no way to know the file
was not altered, or that it belongs to the person presenting it, or when it was
recorded and whether it is still fresh. Second, \emph{control}: once copied, the
user has no mechanism to revoke a baseline that leaked or to signal that it has
been superseded by a recalibration. A public commitment plus a consent envelope
supplies exactly these missing pieces without becoming a data-storage system for
the biometric.

\section{Threat and Trust Model}

We assume: (i) the user's device holds the raw baseline and the signing key and
is trusted; (ii) the ledger is an untrusted-but-tamper-evident public bulletin
board---anything written to it may be read by anyone forever; (iii) recipient
apps are semi-trusted---they should receive only the scope the user approves and
must be able to \emph{verify} what they receive; and (iv) RPC endpoints (RPC
Fast, public devnet) are untrusted for confidentiality but are used only for
integrity-checkable reads and for broadcasting already-signed transactions.

The design goals follow directly:
\begin{itemize}
\item \textbf{Confidentiality:} raw parameters never leave the device except
  under explicit, scoped user consent, and \emph{never} appear on-chain.
\item \textbf{Integrity:} a recipient can detect any modification of the
  baseline it received.
\item \textbf{Provenance \& ownership:} the commitment is bound to the owner's
  wallet and a timestamp.
\item \textbf{Revocability:} the owner can invalidate a baseline.
\item \textbf{Minimality:} the on-chain footprint is a hash and small metadata,
  nothing more.
\end{itemize}

\section{Design}

\subsection{The baseline manifest (on-chain envelope)}

Anchoring a baseline writes a single memo-program transaction whose payload is a
compact JSON envelope. Crucially, the envelope carries a \emph{commitment}, not
the data:

\begin{lstlisting}
{ "type": "baseline_commitment", "v": 1,
  "owner": "<wallet>",
  "signal": "eeg",              // eeg|emg|hrv|...
  "model": "gmm-arousal",
  "modelVersion": "1.0.0",
  "featureSchema": "muse-bandpower-4",
  "featureSchemaHash": "9235...e7d13",
  "baselineHash": "509c...da1ca",   // SHA-256 of params
  "dims": { "features": 4, "clusters": 3 },
  "recordedSeconds": 300,
  "validUntil": "2027-07-31T00:00:00Z",
  "ts": "2026-07-31T00:00:00Z" }
\end{lstlisting}

The \code{baselineHash} is the load-bearing field. \code{dims},
\code{recordedSeconds}, and \code{validUntil} are deliberately coarse:
dimensionality and recording length are context, not identity, and the
recalibration horizon lets a verifier reject a stale baseline. \code{signal} and
\code{model} are short tags so a marketplace of apps can discover which baselines
they can consume.

\subsection{Canonical commitment}

Let $P$ be the fitted parameter object. Define
$\mathrm{canon}(P)$ as the recursively key-sorted, whitespace-free JSON
serialization of $P$, and
\[
\mathrm{baselineHash} \;=\; \mathrm{SHA\text{-}256}\big(\mathrm{canon}(P)\big).
\]
Key sorting makes the commitment independent of how a producer happened to
assemble the object. Because a baseline anchored by one app must verify in
another---possibly written in a different language---$\mathrm{canon}$ also pins
number formatting: whole-valued numbers render without a decimal point
($2.0\!\to\!\texttt{2}$, $0.0\!\to\!\texttt{0}$) and all others via shortest
round-trip, so a Dart (Flutter) producer and a Node.js verifier emit
byte-identical strings. Without this rule the two diverge---Dart's default
encoder renders $0.0$ as \code{"0.0"} where JavaScript renders \code{"0"}---and
the commitments silently mismatch. The same rule is applied to the
feature-schema spec to yield \code{featureSchemaHash}, which lets a consumer
confirm it is extracting features the same way the baseline expects---a
necessary condition for a baseline to transfer correctly.

\subsection{Anchor and verify}

The lifecycle mirrors, and reuses, HealthLog's session-consent machinery:

\begin{lstlisting}
Record baseline (once)
  -> canon + SHA-256  -> baselineHash
    -> assert envelope carries no raw params
      -> memo tx (anchor)  ==> tx signature
--------------------------------------------------
Consumer app receives baseline (out of band)
  -> re-hash received params
    -> read commitment at tx (via RPC Fast)
      -> check: type, hash==, owner==, fresh, no-raw
        -> VERIFIED  => load and use; no re-recording
\end{lstlisting}

Verification is a pure function of the received parameters and the on-chain
commitment; it requires no trusted third party. If the received parameters were
altered by even one value, the recomputed hash diverges and verification fails.

\subsection{Scoped access to the parameters}

Anchoring establishes integrity and ownership; it does not move the baseline.
Sharing the actual parameters is governed by the same scoped, revocable consent
tokens HealthLog already uses for sessions, with a baseline-specific scope
ladder: \emph{metadata only} $\to$ \emph{feature schema only} $\to$ \emph{full
baseline parameters}. The most sensitive scope---handing over the fitted
model---is opt-in and time-boxed, and the on-chain revocation makes a leaked
grant detectable.

\subsection{RPC rails: RPC Fast for verification}

The value of anchoring is realized at \emph{verification} time, when an app that
has never met the user reads the commitment and decides whether to trust a
baseline. That read must be fast and reliable, which is why RPC Fast---a
low-latency Solana RPC provider---serves as the production verification rail on
mainnet. Writing an anchor costs a single memo transaction
($\approx\!5000$ lamports, sub-cent); reading a commitment is a single
\code{getTransaction}. Our free demonstration anchors on devnet (public faucet)
so the full mechanism is exercised end-to-end at zero cost; the anchor and
verify logic are byte-for-byte identical across networks.

\section{Implementation}

The implementation has three parts, all in the HealthLog repository.

\paragraph{Flutter app (\code{lib/}).} A \code{RpcConfig} centralizes the RPC
endpoint and is injected at run time via \code{-{}-dart-define}, so an RPC Fast
API key is never committed. New domain models (\code{WearableBaseline},
\code{BaselineManifest}) and a \code{BaselineAnchorService} interface with mock
and devnet implementations mirror the existing session/consent services. A
\emph{Baseline Passport} screen lets a user load demo EEG/EMG/HRV baselines,
anchor them, and run the verification checklist. The commitment reuses the app's
existing canonical-JSON hashing so the Dart and reference implementations share
one rule.

\paragraph{Reference tool (\code{tools/baseline\_anchor/}).} A small,
dependency-light Node.js tool (\code{anchor\_baseline.mjs},
\code{verify\_baseline.mjs}) implements the same canonicalization, the same
envelope, and a runtime privacy gate that refuses to send any envelope
containing a payload-shaped forbidden key (model means, covariances, weights,
scaler statistics, raw beat series, MVC vectors, \ldots) or an over-long numeric
array. It reads RPC Fast credentials from a local, git-ignored config and signs
with a local keypair. This tool produced every result in \S\ref{sec:eval}.

\paragraph{Privacy gate and tests.} Both implementations enforce the invariant
\emph{``the envelope is a commitment plus shape metadata, nothing more''} at
runtime, and an automated test suite asserts that (i) the commitment is
deterministic and key-order independent, (ii) no forbidden key or raw value
appears in the memo text, (iii) a leaky envelope is rejected, and (iv) a
tampered parameter set fails verification.

\section{Evaluation}
\label{sec:eval}

\paragraph{Real anchors.} We anchored one baseline of each modality on Solana
devnet and verified each round-trip. All six checks
(found on-chain; is a baseline commitment; hash matches; owner matches; within
horizon; no raw parameters) passed for every baseline.

\begin{table}[h]
\centering
\footnotesize
\begin{tabular}{@{}llr@{}}
\toprule
Signal & \code{baselineHash} (prefix) & Env.\ bytes \\
\midrule
EEG (gmm-arousal)      & \code{509c8630cca7\ldots} & 476 \\
EMG (mvc-normalization)& \code{2b9f99cd12d7\ldots} & 482 \\
HRV (resting-norms)    & \code{929e9e779c58\ldots} & 486 \\
\bottomrule
\end{tabular}
\caption{Baselines anchored on Solana devnet from owner wallet
\code{EKRgJifhv4xs\ldots s4sg}. Envelopes are $<\!500$ bytes---well within a
single memo transaction. Full signatures are pinned in the repository README and
receipts.}
\label{tab:anchors}
\end{table}

\paragraph{Tamper detection binds.} Flipping a single GMM weight
($0.42\!\to\!0.43$, compensated elsewhere) changed the commitment from
\code{509c8630\ldots} to \code{bf3eee54\ldots}, and verification against the
unchanged on-chain commitment correctly reported \textsc{failed}. Integrity is
not advisory; it is cryptographically bound to every value in the parameter set.

\paragraph{Cross-language parity.} The two independent implementations agree.
For all three sample baselines, the Dart (Flutter) hasher and the Node.js
reference tool produce byte-identical commitments---enforced as an automated
golden test (\code{test/baseline\_parity\_test.dart}) that fails if either
canonical serializer drifts. This is the property that makes ``anchor in app A,
verify in app B'' real rather than aspirational; before the shared number rule,
the same EEG baseline hashed to \code{a05fdf36\ldots} in Dart versus
\code{509c8630\ldots} in Node.

\paragraph{Verification latency.} Over eight \code{getSlot} probes against the
RPC Fast mainnet endpoint we observed 101/106/371\,ms
(min/median/max), mean 162\,ms, dominated by two cold outliers. A full baseline
verification is a single \code{getTransaction} of comparable cost plus local
hashing (microseconds), so a consumer app can gate ``trust this baseline'' on a
sub-second, one-round-trip check.

\paragraph{Reuse without re-recording (the payoff).} A consumer that verifies a
baseline can immediately \emph{use} it. We implement the consumer side as a
self-contained full-covariance GMM evaluator (the same log-Gaussian form the
app's calibration uses) that standardizes a fresh feature vector with the
baseline's own scaler and returns an arousal label. Feature vectors placed at a
cluster's mean classify to that cluster's label (high$\to$high, low$\to$low,
medium$\to$medium; responsibilities sum to one), all under automated test---so a
second app runs inference from a verified baseline it never recorded, which is
the entire point.

\paragraph{Cost.} Anchoring is one memo transaction. On devnet it is free; on
mainnet it is the base fee ($\approx\!5000$ lamports, well under one US cent),
paid once per baseline (or per recalibration). Reads are free RPC calls.

\section{Security and Privacy Analysis}

\paragraph{What is and is not on-chain.} On-chain: a hash, an owner wallet, short
type/model tags, coarse dimensionality, two timestamps. Off-chain and encrypted:
every fitted value, every feature vector, every raw sample. A hash is a one-way
commitment; recovering parameters from \code{baselineHash} is infeasible.

\paragraph{Biometric-template leakage under sharing.} The commitment protects the
data \emph{on the ledger}, but the sensitive scope of the consent flow does hand
the fitted parameters to a recipient. Fitted model parameters can encode
identity-linked structure; sharing them is a genuine biometric-disclosure
decision, not a metadata one. Our design makes that decision explicit,
time-boxed, and revocable, but does not yet make the shared template
\emph{cancelable}. Cryptographic template protection (\S\ref{sec:future}) is the
right long-term answer; until then the honest framing is: anchoring removes the
\emph{integrity/provenance} risk of portability, while parameter sharing remains
a consented disclosure.

\paragraph{Linkage and correlation.} Multiple commitments from one wallet are
publicly linkable, and \code{signal}/\code{dims}/\code{recordedSeconds} leak
coarse facts (``this wallet has an EEG and an HRV baseline''). A privacy-seeking
user can anchor under rotating or per-baseline wallets at the cost of a unified
passport. We treat wallet hygiene as a user-facing choice, surfaced but not
mandated.

\paragraph{Revocation semantics.} Revocation is a second on-chain event
referencing the baseline; a correct verifier treats a baseline whose revocation
it can see as untrusted. Because the ledger is append-only, revocation cannot
retract data already shared out of band---it can only stop \emph{future} trust.
This is the standard limitation of public-ledger consent and we state it plainly.

\paragraph{Freshness and drift.} Physiology drifts: an HRV resting baseline ages
in weeks, an EEG arousal model more slowly. \code{validUntil} encodes a
recalibration horizon and verification rejects a stale commitment, converting a
silent correctness bug (using an outdated baseline) into an explicit, checkable
failure.

\paragraph{Key management.} The owner wallet is the root of authority. In the
demo it is a device-held keypair; a production deployment should use a hardware
or mobile-wallet signer. Loss of the key means loss of the ability to revoke or
re-anchor, though not disclosure of past data.

\section{Limitations and Honest Positioning}

We deliberately avoid overclaiming.

\begin{itemize}
\item \textbf{Commitment, not proof of quality.} The chain attests that a
  specific parameter set existed and is owned and unmodified. It does \emph{not}
  attest that the baseline was computed correctly, from enough data, or without
  bias. ``This baseline came from $\ge\!5$ minutes of clean EEG'' is a
  zero-knowledge claim we do not yet make (\S\ref{sec:future}).
\item \textbf{Memo now, program state next (written, not yet deployed).} The
  live anchor is an SPL-memo transaction (the MVP). A structured Anchor
  program---per-baseline PDAs keyed by (owner, baseline hash) with native
  revocation---is now written and build-ready (\code{anchor/}) but not yet
  deployed; the memo path remains the tested default. Because the program
  commits the identical hash, the two are interchangeable, and a verifier reads
  a PDA instead of parsing a memo log.
\item \textbf{Cross-language canonicalization (resolved, with a residual edge).}
  The Dart and Node implementations now emit byte-identical canonical strings
  and hashes for realistic baselines, verified by a golden test
  (\S\ref{sec:eval}). The pinned number rule covers normal-range decimals; we do
  not guarantee agreement for pathological magnitudes (values requiring
  exponent notation), which do not arise in EEG/EMG/HRV baselines but which a
  formal spec would still need to fix, e.g.\ by additionally content-addressing
  the transmitted byte blob.
\item \textbf{No cryptographic template protection yet.} As above, the sensitive
  sharing scope discloses a real template.
\item \textbf{Devnet demonstration.} Reported anchors are on devnet; the mainnet
  path is a one-line configuration change and a funded key, and RPC Fast is
  already the live mainnet read rail we measured against.
\end{itemize}

\section{Related Work}

Our design sits at the intersection of self-sovereign identity and biosignal
engineering. \emph{Verifiable credentials} and \emph{decentralized identifiers}
formalize owner-controlled, verifiable claims; we specialize the pattern to a
mutable, drift-bounded biometric baseline and add a recalibration horizon.
\emph{Decentralized health-data} efforts typically anchor document hashes or
access grants; we anchor a \emph{reusable computational artifact}---a
calibration model---and make its \emph{reuse without re-recording} the goal.
From the biosignal side, \emph{transfer learning and cross-session calibration}
for BCIs study how to reduce or transport calibration; our contribution is
orthogonal and complementary: a portable, verifiable container so a transported
baseline can be \emph{trusted} by the receiving system. Finally, \emph{biometric
template protection} (cancelable biometrics, fuzzy commitments) is the body of
work we point to for hardening the parameter-sharing scope.

\section{Future Work}
\label{sec:future}

\begin{itemize}
\item \textbf{Zero-knowledge quality proofs.} Prove ``computed from
  $\ge N$ minutes of artifact-clean signal'' or ``fit converged'' without
  revealing the data or the parameters, upgrading the commitment from
  \emph{it exists} to \emph{it is good}.
\item \textbf{Anchor program with per-baseline accounts.} Native, queryable
  state and revocation flags; a public index of ``baselines this wallet will
  share, and under what scope.''
\item \textbf{Cancelable-biometric templates.} Share a revocable, non-invertible
  transform of the baseline so the sensitive scope no longer discloses a
  reusable template.
\item \textbf{Cross-app SDK.} A drop-in ``verify \& load a HealthLog baseline''
  client so any BCI/wearable app can consume portable baselines, turning the
  passport into an ecosystem rather than a single app's feature.
\end{itemize}

\section{Conclusion}

Baselines are the hidden tax of consumer neurotechnology: costly to collect,
slow to change, and trapped in silos. We showed that a baseline can instead be a
portable, owner-controlled asset by anchoring only its cryptographic commitment
and a thin metadata envelope on a public ledger, keeping the biometric itself
encrypted off-chain, and letting any app verify integrity, provenance, and
freshness before reuse---no re-recording. The mechanism is implemented,
integrated into HealthLog with RPC Fast as the verification rail, and
demonstrated end-to-end on real EEG, EMG, and HRV baselines with cryptographic
tamper detection and sub-second verification. We were equally explicit about what
remains: quality proofs and template protection are future work, and we resist
claiming them prematurely. The result is a concrete step toward wearable
baselines that people \emph{own} and can \emph{carry}.

\begin{thebibliography}{9}
\footnotesize
\bibitem{healthlog} HealthLog Protocol, ``User-owned wearable data with
Solana-based consent and provenance,'' Solana Frontier Hackathon (Colosseum),
2026.
\bibitem{vc} W3C, ``Verifiable Credentials Data Model,'' W3C Recommendation.
\bibitem{did} W3C, ``Decentralized Identifiers (DIDs),'' W3C Recommendation.
\bibitem{bcitl} V.\ Jayaram et al., ``Transfer learning in brain--computer
interfaces,'' \emph{IEEE Computational Intelligence Magazine}, 2016.
\bibitem{cancelable} N.\ Ratha, J.\ Connell, R.\ Bolle, ``Enhancing security and
privacy in biometrics-based authentication systems,'' \emph{IBM Systems
Journal}, 2001.
\bibitem{hrv} Task Force of ESC/NASPE, ``Heart rate variability: standards of
measurement, physiological interpretation, and clinical use,'' \emph{Circulation},
1996.
\bibitem{solana} A.\ Yakovenko, ``Solana: A new architecture for a high
performance blockchain,'' whitepaper, 2018.
\end{thebibliography}

\end{document}
